% REPORT-SHARD: 02_states_and_equivariance
\section{Wilson's triangle, effective states, and symmetry reduction}

\subsection{The two directions of the construction}

At fixed coarse scale $k$, the Wilson row consists of actual microscopic
states pulled back from finer systems:
\[
 \omega_{k,n}=\omega_n\circ\alpha_n^k,\qquad n\ge k.
\]
The horizontal limit $n\to\infty$ defines $\omega_k$; compatibility of these
limits then defines a state on the observable inductive limit. The diagram
separates state convergence from merely comparing operators at two scales:
\[
\begin{array}{ccccccc}
 &&&&\omega_{k+2,k+2}&\longrightarrow\cdots&\omega_{k+2}\\
 &&&&\downarrow&&\downarrow\\
 &&\omega_{k+1,k+1}&\longrightarrow&\omega_{k+1,k+2}&\longrightarrow\cdots&\omega_{k+1}\\
 &&\downarrow&&\downarrow&&\downarrow\\
 \omega_{k,k}&\longrightarrow&\omega_{k,k+1}&\longrightarrow&\omega_{k,k+2}
 &\longrightarrow\cdots&\omega_k
\end{array}
\]
Each horizontal row consists of functionals on its fixed coarse algebra.
Vertical arrows are pullbacks by the appropriate refinement map. The
rightmost column is the compatible limiting family.

WT-1 above gives norm convergence with explicit tails. It is stronger than
the existence of an unspecified subsequential scaling limit. WT-2 turns the
tail into an algorithm for finite-observable expectations only when a
computable modulus and a computable upper bound on the observable norm are
supplied. These are analytic inputs to be produced by a concrete backend.

\subsection{A finite Gram consequence and a necessary warning}

For $a_1,\ldots,a_p\in A_k$, form
$G_n(i,j)=\omega_{k,n}(a_i^*a_j)$ and
$G(i,j)=\omega_k(a_i^*a_j)$. WT-3 gives
\[
 \|G_n-G\|\le t_k(n)\sum_i\|a_i\|^2.
\]
For a coefficient vector $z$, the quadratic form difference is the
state difference on $b^*b$, $b=\sum_i z_i a_i$. WT-1 bounds it by
$t_k(n)\|b\|^2$; Cauchy--Schwarz bounds
$\|b\|^2\le\|z\|_2^2\sum_i\|a_i\|^2$. Taking the supremum proves the claim.

Gram convergence alone does not control kernels: positive matrices
$\operatorname{diag}(1,1/n)$ converge to $\operatorname{diag}(1,0)$.
Likewise $R\Omega=0$ does not imply $R=0$ on the observable orbit.
In two dimensions, choose $R=|1\rangle\langle1|$, $\Omega=|0\rangle$
and $B=|1\rangle\langle0|+|0\rangle\langle1|$. Then $R\Omega=0$ but
$RB\Omega=|1\rangle$. This is the explicit VR-2 counterexample.

\sources{\path{cft_machine/research/wilson_triangle.md}, WT-1--WT-3;
\path{cft_machine/research/virasoro_routes.md}, VR-2.}

\subsection{Finite-group fixed points commute with a strict algebra limit}

Suppose a fixed finite group $G$ acts by automorphisms on every $A_k$,
and all unital injective refinements are equivariant. Then
\[
 \varinjlim A_k^G\ \cong\ (\varinjlim A_k)^G.
\]
One inclusion is immediate. For the other, approximate any invariant element
of the limit by a local element and average that approximant over $G$.
The average stays local and invariant, and its error is no larger because
the finite-group conditional expectation is contractive.

This is a strict C*-algebra statement. It does not automatically transport
local nets, state-dependent weak closures, twisted sectors, or extensions.
Those distinctions were retained in the broader orbifold proposal.

\subsection{Invariant states and their GNS representation}

EQ-1 adds invariant microscopic states. Equivariance makes every Wilson row
invariant, and restricting it to fixed points cannot increase its dual-norm
drift. In the full limiting GNS representation,
\[
 U_g[a]=[\beta_g(a)],\qquad
 P_G=\frac1{|G|}\sum_gU_g
\]
are well-defined unitaries and the projection onto invariant vectors.
The identity
$P_G\pi(a)\Omega=\pi(|G|^{-1}\sum_g\beta_g(a))\Omega$ shows
\[
 \overline{\pi(A^G)\Omega}=\mathcal H^G.
\]
Thus the fixed-algebra vacuum GNS is exactly the invariant cyclic subspace.

EQ-2 proves the represented closure identity on the full Hilbert space,
\[
 \pi(A^G)''=(\pi(A)'')^G.
\]
For an invariant weak-limit operator, average a weakly converging net of
represented algebra elements. Finite averaging is weak-operator continuous,
which gives the reverse inclusion. The equality on the full Hilbert space
must be distinguished from restricting to the vacuum space $\mathcal H^G$.

Local versions require the action and the continuum identification to
preserve the relevant interval algebras. This is exactly why independent
factor parities are appropriate for products of Ising nets, while an arbitrary
orbifold construction still needs additional local and sector arguments.

\sources{\path{cft_machine/research/equivariant_oar.md}, EQ-1/EQ-2;
\path{cft_machine/research/realisation_and_fidelity.md}, fixed-point lemmas;
Fewster--Rejzner source lines 1255--1268 for represented weak closures.}
